% Generated draft: requires mathematical and source review.
\documentclass[11pt]{article}
\usepackage[margin=25mm]{geometry}
\usepackage{fontspec}
\setmainfont{Latin Modern Roman}
\usepackage{amsmath,amssymb,mathtools,unicode-math}
\setmathfont{Latin Modern Math}
\usepackage{enumitem,needspace,xurl,hyperref}
\setlength{\parindent}{0pt}
\setlength{\parskip}{5pt}
\setlength{\emergencystretch}{3em}
\newcommand{\N}{\mathbb{N}}
\newcommand{\Z}{\mathbb{Z}}
\newcommand{\Q}{\mathbb{Q}}
\newcommand{\R}{\mathbb{R}}
\newcommand{\C}{\mathbb{C}}
\newcommand{\F}{\mathbb{F}}
\newcommand{\A}{\mathbb{A}}
\newcommand{\PP}{\mathbb{P}}
\newcommand{\E}{\mathbb{E}}
\newcommand{\eps}{\varepsilon}
\newcommand{\dd}{\,\mathrm{d}}
\begin{document}
\section*{UnsolvedMath ID: 30004179}
\textbf{Source problem code:} catalogue entry
\subsection{Definitions and mathematical statement}
A surface is a smooth manifold of dimension two; the supplied record does not specify compactness, connectedness, orientability, or whether boundary is allowed. A surface diffeomorphism is an invertible differentiable map \(f\colon S\to S\) with differentiable inverse, with the precise differentiability class left unspecified by the source.

For a specified class \(\mathcal D(S)\) of surface diffeomorphisms and a specified convention for topological entropy, write \(h_{\mathrm{top}}(f)\) for the topological entropy of \(f\), and set
\[
\mathcal Z(S)=\{f\in\mathcal D(S):h_{\mathrm{top}}(f)=0\},
\qquad
\mathcal P(S)=\{f\in\mathcal D(S):h_{\mathrm{top}}(f)>0\}.
\]
For compact surfaces, topological entropy has its usual meaning for continuous self-maps of compact spaces. If noncompact surfaces are intended, an entropy convention must be specified; compactness is not an assumption supplied by this record.

\textbf{Research problem.} Describe and classify the dynamics of the diffeomorphisms in \(\mathcal Z(S)\), and describe the mechanisms by which zero-topological-entropy surface diffeomorphisms bifurcate to dynamics in \(\mathcal P(S)\).

\textbf{Unresolved formulation qualifications.} This is a broad research request rather than a uniquely specified theorem or conjecture. The source does not fix the class of surfaces, differentiability class, equivalence relation or invariants for classification, topology on the space of maps, or the families and parameter conditions used to define bifurcation. It does not specify whether the requested mechanisms concern all zero-entropy maps, generic maps, particular families, or particular subclasses. These choices and the associated quantifiers must be supplied before a precise classification or bifurcation assertion can be stated.
\subsection{Short English statement}
Describe the possible dynamics of surface diffeomorphisms with zero topological entropy and explain how positive entropy can arise through bifurcation. The supplied question leaves the surface class, regularity, classification criterion, and meaning of bifurcation unspecified, so it determines a research program rather than a precise conjecture.
\subsection{Sources}
\begin{itemize}
\item Local UnsolvedMath registry: permanent ID \(30004179\), title ``Zero-Entropy Surface Dynamics and Their Bifurcations,'' source problem code \texttt{catalogue entry}, in the Oberwolfach Reports Open Problems collection. The supplied proposed year is 2019; no individual proposer is identified.
\item Supplied report citation: \emph{Dynamische Systeme (2020)}, Oberwolfach Reports, DOI \texttt{10.4171/owr/2019/31}, pp.\ 1855--1917. \url{https://doi.org/10.4171/owr/2019/31}. The background describes these reports as records of mathematical workshops and their open-problem discussions. No page or speaker for this individual question is supplied.
\item Supplied literature reference: ``Crovisier--Pujals--Tresser, zero entropy.'' \url{https://arxiv.org/abs/2005.14278}. No specific theorem or passage supporting the extracted question or assessment is identified.
\item Supplied literature reference: ``original OWR report.'' \url{https://doi.org/10.4171/owr/2019/31}.
\item Literature assessment dated exactly \texttt{2026-08-22}: the supplied curation labels the problem open and states that available theorems cover restricted classes without settling the general classification and bifurcation questions. This is supplied dated curation derived from the OWR material, not an independently verified literature assessment; time-sensitive status and cited claims require independent verification.
\item The registry labels the statement \texttt{corrected\_verified} and says its mathematical content was preserved while notation, typography, and sentence structure were repaired for publication. This is supplied editorial metadata, not independent mathematical verification of this draft.
\end{itemize}
\end{document}
